
\documentclass[12pt,a4paper]{article}
\usepackage[UTF8]{ctex}
\usepackage{geometry}
\usepackage{amsmath,amssymb}
\usepackage{listings}
\usepackage{xcolor}
\usepackage{longtable}
\usepackage{hyperref}
\geometry{margin=1in}

\lstset{
    language=C++,
    basicstyle=\ttfamily\small,
    keywordstyle=\color{blue},
    commentstyle=\color{green!50!black},
    stringstyle=\color{red!60!black},
    breaklines=true,
    showstringspaces=false,
    frame=single
}

\title{KMP 知识讲解与 OJ 练习题}
\author{pinqu（AI 辅助）}
\date{2026.4.13}

\begin{document}
\maketitle

\section{KMP 基础知识讲解}

\subsection{KMP 要解决什么问题}
KMP（Knuth-Morris-Pratt）算法用于解决字符串匹配问题：给定主串 $S$ 和模式串 $P$，判断 $P$ 是否在 $S$ 中出现，或者找出它出现的位置。

暴力匹配(BF算法)在失配时会让主串指针回退，最坏时间复杂度为 $O(nm)$。  
KMP 的核心思想是：\textbf{主串指针不回退，只回退模式串指针}，从而把复杂度优化到 $O(n+m)$。

\subsection{next 数组的意义}
当模式串在位置 $j$ 失配时，KMP 不会从头开始比较，而是跳到 $next[j]$ 对应的位置继续匹配。

它本质上利用了模式串自身的前后缀信息。  
例如字符串 \texttt{ababaca}，在某个位置失配时，如果已知前面已经匹配了一段，那么模式串内部已经包含的信息可以帮助我们跳过重复比较。

\subsection{两种常见写法}
常见有两种体系：

\begin{itemize}
    \item \texttt{next[0] = -1} 写法
    \item \texttt{lps[0] = 0} 写法
\end{itemize}

本文统一使用 \texttt{next[0] = -1} 的写法。

\subsection{next 数组构造模板}
\begin{lstlisting}
vector<int> getNext(const string& p) {
    int m = p.size();
    vector<int> nxt(m, -1);

    int i = 0, j = -1;
    while (i < m - 1) {
        if (j == -1 || p[i] == p[j]) {
            i++;
            j++;
            nxt[i] = j;
        } else {
            j = nxt[j];
        }
    }
    return nxt;
}
\end{lstlisting}

\subsection{KMP 匹配模板}
\begin{lstlisting}
int kmp(const string& s, const string& p) {
    if (p.empty()) return 0;

    vector<int> nxt = getNext(p);
    int i = 0, j = 0;
    int n = s.size(), m = p.size();

    while (i < n && j < m) {
        if (j == -1 || s[i] == p[j]) {
            i++;
            j++;
        } else {
            j = nxt[j];
        }
    }

    if (j == m) return i - j;
    return -1;
}
\end{lstlisting}

\subsection{最容易出错的点}
\begin{enumerate}
    \item 判断顺序必须写成 \texttt{j == -1 || s[i] == p[j]}，否则可能访问 \texttt{p[-1]}。
    \item 构造 \texttt{next} 时循环通常写成 \texttt{while (i < m - 1)}。
    \item 失配时回退的是 \texttt{j}，不是 \texttt{i}。
    \item 不要把 \texttt{next[0] = -1} 和 \texttt{lps[0] = 0} 两套写法混在一起。
\end{enumerate}

\section{OJ 题目与参考代码}

\subsection*{题目 1：字符串首次匹配}

\textbf{题目描述}  
给定两个字符串 $S$ 和 $P$，求模式串 $P$ 在主串 $S$ 中第一次出现的位置。若不存在，输出 $-1$。

\textbf{输入格式}
\begin{itemize}
    \item 第一行输入字符串 $S$
    \item 第二行输入字符串 $P$
\end{itemize}

\textbf{输出格式}  
输出一个整数。

\textbf{数据范围}  
$1 \le |S|, |P| \le 10^5$

\textbf{样例输入}
\begin{lstlisting}
helloabcworld
abc
\end{lstlisting}

\textbf{样例输出}
\begin{lstlisting}
5
\end{lstlisting}

\textbf{参考代码}
\begin{lstlisting}
#include <iostream>
#include <vector>
#include <string>
using namespace std;

vector<int> getNext(const string& p) {
    int m = p.size();
    vector<int> nxt(m, -1);
    int i = 0, j = -1;
    while (i < m - 1) {
        if (j == -1 || p[i] == p[j]) {
            ++i; ++j;
            nxt[i] = j;
        } else {
            j = nxt[j];
        }
    }
    return nxt;
}

int kmp(const string& s, const string& p) {
    if (p.empty()) return 0;
    vector<int> nxt = getNext(p);
    int i = 0, j = 0;
    while (i < (int)s.size() && j < (int)p.size()) {
        if (j == -1 || s[i] == p[j]) {
            ++i; ++j;
        } else {
            j = nxt[j];
        }
    }
    return (j == (int)p.size()) ? i - j : -1;
}

int main() {
    string s, p;
    cin >> s >> p;
    cout << kmp(s, p) << "\n";
    return 0;
}
\end{lstlisting}

\subsection*{题目 2：输出 next 数组}

\textbf{题目描述}  
给定一个模式串 $P$，请输出它的 \texttt{next} 数组（采用 \texttt{next[0] = -1} 写法）。

\textbf{输入格式}  
输入一个字符串 $P$。

\textbf{输出格式}  
输出一行 $n$ 个整数，表示 \texttt{next} 数组。

\textbf{样例输入}
\begin{lstlisting}
ababaca
\end{lstlisting}

\textbf{样例输出}
\begin{lstlisting}
-1 0 0 1 2 3 0
\end{lstlisting}

\textbf{参考代码}
\begin{lstlisting}
#include <iostream>
#include <vector>
#include <string>
using namespace std;

vector<int> getNext(const string& p) {
    int m = p.size();
    vector<int> nxt(m, -1);
    int i = 0, j = -1;
    while (i < m - 1) {
        if (j == -1 || p[i] == p[j]) {
            ++i; ++j;
            nxt[i] = j;
        } else {
            j = nxt[j];
        }
    }
    return nxt;
}

int main() {
    string p;
    cin >> p;
    vector<int> nxt = getNext(p);
    for (int i = 0; i < (int)nxt.size(); ++i) {
        cout << nxt[i] << (i + 1 == (int)nxt.size() ? '\n' : ' ');
    }
    return 0;
}
\end{lstlisting}

\subsection*{题目 3：找出所有匹配位置}

\textbf{题目描述}  
给定主串 $S$ 和模式串 $P$，输出模式串在主串中出现的所有位置，允许重叠匹配。

\textbf{输入格式}
\begin{itemize}
    \item 第一行输入字符串 $S$
    \item 第二行输入字符串 $P$
\end{itemize}

\textbf{输出格式}  
按升序输出所有匹配起始位置，空格分隔。若不存在，输出 \texttt{-1}。

\textbf{样例输入}
\begin{lstlisting}
aaaaa
aa
\end{lstlisting}

\textbf{样例输出}
\begin{lstlisting}
0 1 2 3
\end{lstlisting}

\textbf{参考代码}
\begin{lstlisting}
#include <iostream>
#include <vector>
#include <string>
using namespace std;

vector<int> getNext(const string& p) {
    int m = p.size();
    vector<int> nxt(m, -1);
    int i = 0, j = -1;
    while (i < m - 1) {
        if (j == -1 || p[i] == p[j]) {
            ++i; ++j;
            nxt[i] = j;
        } else {
            j = nxt[j];
        }
    }
    return nxt;
}

int main() {
    string s, p;
    cin >> s >> p;
    vector<int> nxt = getNext(p);
    vector<int> ans;

    int i = 0, j = 0;
    while (i < (int)s.size()) {
        if (j == -1 || s[i] == p[j]) {
            ++i; ++j;
        } else {
            j = nxt[j];
        }

        if (j == (int)p.size()) {
            ans.push_back(i - j);
            j = nxt[j - 1] + 1;
        }
    }

    if (ans.empty()) {
        cout << -1 << "\n";
    } else {
        for (int k = 0; k < (int)ans.size(); ++k) {
            cout << ans[k] << (k + 1 == (int)ans.size() ? '\n' : ' ');
        }
    }
    return 0;
}
\end{lstlisting}

\subsection*{题目 4：统计模式串出现次数}

\textbf{题目描述}  
给定字符串 $S$ 和模式串 $P$，统计 $P$ 在 $S$ 中出现的次数，允许重叠。

\textbf{输入格式}
\begin{itemize}
    \item 第一行输入字符串 $S$
    \item 第二行输入字符串 $P$
\end{itemize}

\textbf{输出格式}  
输出一个整数。

\textbf{样例输入}
\begin{lstlisting}
abababab
abab
\end{lstlisting}

\textbf{样例输出}
\begin{lstlisting}
3
\end{lstlisting}

\textbf{参考代码}
\begin{lstlisting}
#include <iostream>
#include <vector>
#include <string>
using namespace std;

vector<int> getNext(const string& p) {
    int m = p.size();
    vector<int> nxt(m, -1);
    int i = 0, j = -1;
    while (i < m - 1) {
        if (j == -1 || p[i] == p[j]) {
            ++i; ++j;
            nxt[i] = j;
        } else {
            j = nxt[j];
        }
    }
    return nxt;
}

int main() {
    string s, p;
    cin >> s >> p;
    vector<int> nxt = getNext(p);

    int cnt = 0;
    int i = 0, j = 0;
    while (i < (int)s.size()) {
        if (j == -1 || s[i] == p[j]) {
            ++i; ++j;
        } else {
            j = nxt[j];
        }

        if (j == (int)p.size()) {
            ++cnt;
            j = nxt[j - 1] + 1;
        }
    }

    cout << cnt << "\n";
    return 0;
}
\end{lstlisting}

\subsection*{题目 5：判断重复子串构成}

\textbf{题目描述}  
给定一个字符串 $s$，判断它是否可以由某个更短的子串重复若干次构成。

\textbf{输入格式}  
输入一个字符串 $s$。

\textbf{输出格式}  
若可以，输出 \texttt{true}；否则输出 \texttt{false}。

\textbf{样例输入}
\begin{lstlisting}
ababab
\end{lstlisting}

\textbf{样例输出}
\begin{lstlisting}
true
\end{lstlisting}

\textbf{参考代码}
\begin{lstlisting}
#include <iostream>
#include <vector>
#include <string>
using namespace std;

// 这里用 lps 写法更直观
vector<int> getLPS(const string& s) {
    int n = s.size();
    vector<int> lps(n, 0);
    for (int i = 1, len = 0; i < n; ) {
        if (s[i] == s[len]) {
            lps[i++] = ++len;
        } else if (len > 0) {
            len = lps[len - 1];
        } else {
            lps[i++] = 0;
        }
    }
    return lps;
}

int main() {
    string s;
    cin >> s;
    vector<int> lps = getLPS(s);
    int n = s.size();
    int len = lps[n - 1];
    int cycle = n - len;

    if (len > 0 && n % cycle == 0) cout << "true\n";
    else cout << "false\n";
    return 0;
}
\end{lstlisting}

\subsection*{题目 6：最小循环节长度}

\textbf{题目描述}  
给定一个字符串 $s$，求它的最小循环节长度。若不存在严格循环结构，则输出字符串长度。

\textbf{输入格式}  
输入一个字符串 $s$。

\textbf{输出格式}  
输出一个整数。

\textbf{样例输入}
\begin{lstlisting}
abcabcabcabc
\end{lstlisting}

\textbf{样例输出}
\begin{lstlisting}
3
\end{lstlisting}

\textbf{参考代码}
\begin{lstlisting}
#include <iostream>
#include <vector>
#include <string>
using namespace std;

vector<int> getLPS(const string& s) {
    int n = s.size();
    vector<int> lps(n, 0);
    for (int i = 1, len = 0; i < n; ) {
        if (s[i] == s[len]) {
            lps[i++] = ++len;
        } else if (len > 0) {
            len = lps[len - 1];
        } else {
            lps[i++] = 0;
        }
    }
    return lps;
}

int main() {
    string s;
    cin >> s;
    int n = s.size();
    vector<int> lps = getLPS(s);
    int len = lps[n - 1];
    int cycle = n - len;

    if (len > 0 && n % cycle == 0) cout << cycle << "\n";
    else cout << n << "\n";
    return 0;
}
\end{lstlisting}

\subsection*{题目 7：字符串旋转判断}

\textbf{题目描述}  
给定两个字符串 $s_1$ 和 $s_2$，判断 $s_2$ 是否可以由 $s_1$ 旋转得到。

\textbf{输入格式}
\begin{itemize}
    \item 第一行输入字符串 $s_1$
    \item 第二行输入字符串 $s_2$
\end{itemize}

\textbf{输出格式}  
若可以输出 \texttt{true}，否则输出 \texttt{false}。

\textbf{样例输入}
\begin{lstlisting}
abcde
cdeab
\end{lstlisting}

\textbf{样例输出}
\begin{lstlisting}
true
\end{lstlisting}

\textbf{参考代码}
\begin{lstlisting}
#include <iostream>
#include <vector>
#include <string>
using namespace std;

vector<int> getNext(const string& p) {
    int m = p.size();
    vector<int> nxt(m, -1);
    int i = 0, j = -1;
    while (i < m - 1) {
        if (j == -1 || p[i] == p[j]) {
            ++i; ++j;
            nxt[i] = j;
        } else {
            j = nxt[j];
        }
    }
    return nxt;
}

int kmp(const string& s, const string& p) {
    if (p.empty()) return 0;
    vector<int> nxt = getNext(p);
    int i = 0, j = 0;
    while (i < (int)s.size() && j < (int)p.size()) {
        if (j == -1 || s[i] == p[j]) {
            ++i; ++j;
        } else {
            j = nxt[j];
        }
    }
    return (j == (int)p.size()) ? i - j : -1;
}

int main() {
    string s1, s2;
    cin >> s1 >> s2;
    if (s1.size() != s2.size()) {
        cout << "false\n";
        return 0;
    }
    string doubled = s1 + s1;
    cout << (kmp(doubled, s2) != -1 ? "true" : "false") << "\n";
    return 0;
}
\end{lstlisting}

\subsection*{题目 8：最短回文串（前缀添加）}

\textbf{题目描述}  
给定一个字符串 $s$，你可以在它前面添加若干字符，使整个字符串变成回文串。请输出构造出的最短回文串。

\textbf{输入格式}  
输入一个字符串 $s$。

\textbf{输出格式}  
输出最短回文串。

\textbf{样例输入}
\begin{lstlisting}
abcd
\end{lstlisting}

\textbf{样例输出}
\begin{lstlisting}
dcbabcd
\end{lstlisting}

\textbf{参考代码}
\begin{lstlisting}
#include <iostream>
#include <vector>
#include <string>
#include <algorithm>
using namespace std;

vector<int> getLPS(const string& s) {
    int n = s.size();
    vector<int> lps(n, 0);
    for (int i = 1, len = 0; i < n; ) {
        if (s[i] == s[len]) {
            lps[i++] = ++len;
        } else if (len > 0) {
            len = lps[len - 1];
        } else {
            lps[i++] = 0;
        }
    }
    return lps;
}

int main() {
    string s;
    cin >> s;
    string rev = s;
    reverse(rev.begin(), rev.end());

    string t = s + "#" + rev;
    vector<int> lps = getLPS(t);
    int longest = lps.back();

    string add = s.substr(longest);
    reverse(add.begin(), add.end());
    cout << add + s << "\n";
    return 0;
}
\end{lstlisting}

\section{练习建议（AI assist)}
建议按以下顺序完成练习：

\begin{enumerate}
    \item 先做题 2，自己打印 \texttt{next} 数组。
    \item 再做题 1，练熟 KMP 主流程。
    \item 然后做题 3 和题 4，理解重叠匹配。
    \item 接着做题 5 和题 6，掌握字符串周期。
    \item 最后做题 7 和题 8，理解 KMP 的实际应用。
\end{enumerate}

\end{document}
